% Options for packages loaded elsewhere
% Options for packages loaded elsewhere
\PassOptionsToPackage{unicode}{hyperref}
\PassOptionsToPackage{hyphens}{url}
\PassOptionsToPackage{dvipsnames,svgnames,x11names}{xcolor}
%
\documentclass[
  spanish,
  10pt,
  letterpaper,
  twocolumn]{article}
\usepackage{xcolor}
\usepackage[top=1.0cm,bottom=1.5cm,left=1.5cm,right=1.5cm]{geometry}
\usepackage{amsmath,amssymb}
\setcounter{secnumdepth}{-\maxdimen} % remove section numbering
\usepackage{iftex}
\ifPDFTeX
  \usepackage[T1]{fontenc}
  \usepackage[utf8]{inputenc}
  \usepackage{textcomp} % provide euro and other symbols
\else % if luatex or xetex
  \usepackage{unicode-math} % this also loads fontspec
  \defaultfontfeatures{Scale=MatchLowercase}
  \defaultfontfeatures[\rmfamily]{Ligatures=TeX,Scale=1}
\fi
\usepackage{lmodern}
\ifPDFTeX\else
  % xetex/luatex font selection
  \setmainfont[]{Fira Sans}
\fi
% Use upquote if available, for straight quotes in verbatim environments
\IfFileExists{upquote.sty}{\usepackage{upquote}}{}
\IfFileExists{microtype.sty}{% use microtype if available
  \usepackage[]{microtype}
  \UseMicrotypeSet[protrusion]{basicmath} % disable protrusion for tt fonts
}{}
\usepackage{setspace}
\makeatletter
\@ifundefined{KOMAClassName}{% if non-KOMA class
  \IfFileExists{parskip.sty}{%
    \usepackage{parskip}
  }{% else
    \setlength{\parindent}{0pt}
    \setlength{\parskip}{6pt plus 2pt minus 1pt}}
}{% if KOMA class
  \KOMAoptions{parskip=half}}
\makeatother
% Make \paragraph and \subparagraph free-standing
\makeatletter
\ifx\paragraph\undefined\else
  \let\oldparagraph\paragraph
  \renewcommand{\paragraph}{
    \@ifstar
      \xxxParagraphStar
      \xxxParagraphNoStar
  }
  \newcommand{\xxxParagraphStar}[1]{\oldparagraph*{#1}\mbox{}}
  \newcommand{\xxxParagraphNoStar}[1]{\oldparagraph{#1}\mbox{}}
\fi
\ifx\subparagraph\undefined\else
  \let\oldsubparagraph\subparagraph
  \renewcommand{\subparagraph}{
    \@ifstar
      \xxxSubParagraphStar
      \xxxSubParagraphNoStar
  }
  \newcommand{\xxxSubParagraphStar}[1]{\oldsubparagraph*{#1}\mbox{}}
  \newcommand{\xxxSubParagraphNoStar}[1]{\oldsubparagraph{#1}\mbox{}}
\fi
\makeatother


\usepackage{longtable,booktabs,array}
\newcounter{none} % for unnumbered tables
\usepackage{calc} % for calculating minipage widths
% Correct order of tables after \paragraph or \subparagraph
\usepackage{etoolbox}
\makeatletter
\patchcmd\longtable{\par}{\if@noskipsec\mbox{}\fi\par}{}{}
\makeatother
% Allow footnotes in longtable head/foot
\IfFileExists{footnotehyper.sty}{\usepackage{footnotehyper}}{\usepackage{footnote}}
\makesavenoteenv{longtable}
\usepackage{graphicx}
\makeatletter
\newsavebox\pandoc@box
\newcommand*\pandocbounded[1]{% scales image to fit in text height/width
  \sbox\pandoc@box{#1}%
  \Gscale@div\@tempa{\textheight}{\dimexpr\ht\pandoc@box+\dp\pandoc@box\relax}%
  \Gscale@div\@tempb{\linewidth}{\wd\pandoc@box}%
  \ifdim\@tempb\p@<\@tempa\p@\let\@tempa\@tempb\fi% select the smaller of both
  \ifdim\@tempa\p@<\p@\scalebox{\@tempa}{\usebox\pandoc@box}%
  \else\usebox{\pandoc@box}%
  \fi%
}
% Set default figure placement to htbp
\def\fps@figure{htbp}
\makeatother



\ifLuaTeX
\usepackage[bidi=basic,shorthands=off]{babel}
\else
\usepackage[bidi=default,shorthands=off]{babel}
\fi
\ifPDFTeX
\else
\babelfont{rm}[]{Fira Sans}
\fi
\ifLuaTeX
  \usepackage{selnolig} % disable illegal ligatures
\fi


\setlength{\emergencystretch}{3em} % prevent overfull lines

\providecommand{\tightlist}{%
  \setlength{\itemsep}{0pt}\setlength{\parskip}{0pt}}



 


\usepackage{tikz}
\usepackage{graphicx}
\usepackage{amsmath}
\renewcommand{\arraystretch}{1.4}
\makeatletter
\@ifpackageloaded{caption}{}{\usepackage{caption}}
\AtBeginDocument{%
\ifdefined\contentsname
  \renewcommand*\contentsname{Tabla de contenidos}
\else
  \newcommand\contentsname{Tabla de contenidos}
\fi
\ifdefined\listfigurename
  \renewcommand*\listfigurename{Listado de Figuras}
\else
  \newcommand\listfigurename{Listado de Figuras}
\fi
\ifdefined\listtablename
  \renewcommand*\listtablename{Listado de Tablas}
\else
  \newcommand\listtablename{Listado de Tablas}
\fi
\ifdefined\figurename
  \renewcommand*\figurename{Figura}
\else
  \newcommand\figurename{Figura}
\fi
\ifdefined\tablename
  \renewcommand*\tablename{Tabla}
\else
  \newcommand\tablename{Tabla}
\fi
}
\@ifpackageloaded{float}{}{\usepackage{float}}
\floatstyle{ruled}
\@ifundefined{c@chapter}{\newfloat{codelisting}{h}{lop}}{\newfloat{codelisting}{h}{lop}[chapter]}
\floatname{codelisting}{Listado}
\newcommand*\listoflistings{\listof{codelisting}{Listado de Listados}}
\makeatother
\makeatletter
\makeatother
\makeatletter
\@ifpackageloaded{caption}{}{\usepackage{caption}}
\@ifpackageloaded{subcaption}{}{\usepackage{subcaption}}
\makeatother
\usepackage{bookmark}
\IfFileExists{xurl.sty}{\usepackage{xurl}}{} % add URL line breaks if available
\urlstyle{same}
% fallback for those not using the hyperref driver hyperxmp:
\makeatletter
\@ifundefined{xmpquote}{\newcommand{\xmpquote}[1]{#1}}{}
\makeatother
\hypersetup{
  pdftitle={Práctica -- Fisiología de la Germinación y Estrés Abiótico},
  pdflang={es},
  colorlinks=true,
  linkcolor={blue},
  filecolor={Maroon},
  citecolor={Blue},
  urlcolor={Blue},
  pdfcreator={LaTeX via pandoc}}


\title{\vspace{-1.5em}{\Huge \textbf{Práctica -- Fisiología de la
Germinación y Estrés Abiótico}}\vspace{-1.5em}}
\date{}
\begin{document}
\maketitle


\setstretch{1.2}
\textbf{Fecha}: \hrulefill \hfill \textbf{Grupo}: \hrulefill

\textbf{Nombre:} \hrulefill

\textbf{Objetivos:}

\begin{itemize}
\tightlist
\item
  Evaluar la cinética de germinación y el efecto del potencial osmótico
  y la luz sobre la protrusión radicular.
\item
  Determinar cuantitativamente el Porcentaje Final de Germinación
  (\(\%G\)), el Tiempo Medio de Germinación (\(TMG\)) y la longitud de
  la radícula.
\end{itemize}

\subsection{1. Materiales}\label{materiales}

\begin{itemize}
\tightlist
\item
  Semillas homogéneas (\emph{Phaseolus vulgaris} / frijol o \emph{Zea
  mays} / maíz).
\item
  Cajas de Petri esterilizadas y papel filtro (disco).
\item
  Soluciones de tratamiento: Agua destilada (Control), NaCl 150 mM
  (Estrés Salino) y solución en oscuridad/cubierta de aluminio.
\item
  Regla graduada milimetrada, bisturí y pinzas de precisión.
\item
  Balanza analítica (\(0.001\text{ g}\)).
\item
  Agua destilada, frascos lavadores y papel absorbente.
\item
  Marcador permanente y cinta adhesiva.
\end{itemize}

\subsection{2. Aspectos Teóricos e
Introducción}\label{aspectos-teuxf3ricos-e-introducciuxf3n}

La germinación inicia con la \textbf{embibición} de agua por la semilla
seca y culmina con la protrusión de la radícula a través de las
cubiertas seminales. Este proceso está estrictamente regulado por la
disponibilidad hídrica (\(\Psi_w\)), el balance fitohormonal entre el
\textbf{Ácido Giberélico (\(GA_3\))} (promotor) y el \textbf{Ácido
Abscísico (\(ABA\))} (inhibidor), y señales ambientales como la luz
(fitocromos) y la temperatura.

Defina brevemente:

• \textbf{Fases de la Embibición y Curva Trifásica de Absorción de
Agua:}

\hrulefill

\hrulefill

• \textbf{Efecto del Potencial Osmótico (\(\Psi_s\)) sobre la
Germinación:}

\hrulefill

\hrulefill

• \textbf{Fotoblastismo (Semillas Fotoblásticas Positivas
vs.~Negativas):}

\hrulefill

\hrulefill

\subsection{3. Procedimiento
Experimental}\label{procedimiento-experimental}

\subsubsection{\texorpdfstring{A. Montaje de Tratamientos (\(n=3\) cajas
por
grupo)}{A. Montaje de Tratamientos (n=3 cajas por grupo)}}\label{a.-montaje-de-tratamientos-n3-cajas-por-grupo}

\begin{enumerate}
\def\labelenumi{\arabic{enumi}.}
\tightlist
\item
  Seleccione 225 semillas sanas de tamaño y peso uniforme (25 semillas
  por caja de Petri).
\item
  Prepare 9 cajas de Petri con papel filtro en el fondo y distribuya los
  tratamientos:

  \begin{itemize}
  \tightlist
  \item
    \textbf{T1 (Control - Luz / Agua Destilada):} 3 cajas de Petri
    humedecidas con 5 mL de agua destilada, expuestas a fotoperíodo
    normal.
  \item
    \textbf{T2 (Estrés Salino - NaCl 150 mM):} 3 cajas de Petri
    humedecidas con 5 mL de solución de NaCl 150 mM bajo fotoperíodo
    normal.
  \item
    \textbf{T3 (Oscuridad Continua - Agua Destilada):} 3 cajas de Petri
    humedecidas con 5 mL de agua destilada y envueltas completamente en
    papel aluminio.
  \end{itemize}
\end{enumerate}

\begin{center}
\begin{tikzpicture}[scale=0.78]
  % TRATAMIENTO 1: Control (Luz / Agua)
  \begin{scope}[shift={(0,0)}]
    \draw[thick] (0,0) ellipse (0.75cm and 0.25cm);
    \draw[thick] (-0.75,0) -- (-0.75,0.3) arc (180:360:0.75cm and 0.25cm) -- (0.75,0);
    \fill[gray!20] (-0.73,0.05) arc (180:360:0.73cm and 0.23cm) -- cycle;
    % Semillas germinadas con radícula
    \fill[green!40!black] (-0.3,0.05) circle (1.8pt); \draw[green!60!black, thick] (-0.3,0.05) -- (-0.4,-0.05);
    \fill[green!40!black] (0.1,0.1) circle (1.8pt); \draw[green!60!black, thick] (0.1,0.1) -- (0.2,0.02);
    \fill[green!40!black] (0.3,-0.05) circle (1.8pt); \draw[green!60!black, thick] (0.3,-0.05) -- (0.4,-0.15);
    
    \node[font=\fontsize{6}{7}\bfseries\selectfont, align=center] at (0, -0.5) {T1: Control};
    \node[font=\fontsize{5}{6}\selectfont, gray!50!black, align=center] at (0, -0.8) {Luz / $H_2O$ ($n=3$)};
  \end{scope}

  % TRATAMIENTO 2: Estrés Salino
  \begin{scope}[shift={(2.6,0)}]
    \draw[thick] (0,0) ellipse (0.75cm and 0.25cm);
    \draw[thick] (-0.75,0) -- (-0.75,0.3) arc (180:360:0.75cm and 0.25cm) -- (0.75,0);
    \fill[blue!10] (-0.73,0.05) arc (180:360:0.73cm and 0.23cm) -- cycle;
    % Semillas hinchadas/sin radícula
    \fill[brown!70!black] (-0.3,0.05) circle (2.0pt);
    \fill[brown!70!black] (0.0,0.08) circle (2.0pt);
    \fill[brown!70!black] (0.35,-0.02) circle (2.0pt);
    
    \node[font=\fontsize{6}{7}\bfseries\selectfont, align=center] at (0, -0.5) {T2: Estrés Salino};
    \node[font=\fontsize{5}{6}\selectfont, gray!50!black, align=center] at (0, -0.8) {NaCl 150 mM ($n=3$)};
  \end{scope}

  % TRATAMIENTO 3: Oscuridad
  \begin{scope}[shift={(5.2,0)}]
    \draw[thick, fill=gray!40] (0,0) ellipse (0.75cm and 0.25cm);
    \draw[thick, fill=gray!50] (-0.75,0) -- (-0.75,0.3) arc (180:360:0.75cm and 0.25cm) -- (0.75,0) -- cycle;
    \node[font=\fontsize{5}{6}\bfseries\selectfont, black] at (0,0.15) {ALUMINIO};

    \node[font=\fontsize{6}{7}\bfseries\selectfont, align=center] at (0, -0.5) {T3: Oscuridad};
    \node[font=\fontsize{5}{6}\selectfont, gray!50!black, align=center] at (0, -0.8) {Cubierto ($n=3$)};
  \end{scope}
\end{tikzpicture}
\end{center}

\subsubsection{B. Evaluación Cuantitativa (Día
7)}\label{b.-evaluaciuxf3n-cuantitativa-duxeda-7}

\begin{enumerate}
\def\labelenumi{\arabic{enumi}.}
\tightlist
\item
  A los 7 días de incubación a temperatura constante
  (\(22-25^\circ\text{C}\)), abra las cajas (T3 solo al momento de la
  lectura final).
\item
  Se considera una semilla \textbf{germinada} si la radícula presenta
  una longitud \(\ge 2\text{ mm}\).
\item
  Registre el número total de semillas germinadas por caja (\(N_g\))
  para calcular el Porcentaje de Germinación
  (\(\%G = \frac{N_g}{25} \times 100\)).
\item
  Mida con regla la longitud de la radícula (\(\text{cm}\)) en 5
  semillas germinadas al azar por caja de Petri.
\end{enumerate}

\subsection{4. Registro de Resultados}\label{registro-de-resultados}

\subsubsection{Tabla 1. Porcentaje de Germinación (\%G) e Índice de
Vigor al Día
7}\label{tabla-1.-porcentaje-de-germinaciuxf3n-g-e-uxedndice-de-vigor-al-duxeda-7}

\begin{center}
\resizebox{\linewidth}{!}{%
\begin{tabular}{l|l|c|c|c|c|c}
\hline
\textbf{Tratamiento} & \textbf{Variable} & \textbf{Rép. A} & \textbf{Rép. B} & \textbf{Rép. C} & \textbf{Promedio ($\bar{x}$)} & \textbf{Desv. Est. ($s$)} \\
\hline
1. Control & Germinación ($\%$) & & & & & \\
(Luz / Agua) & Long. Radícula (cm) & & & & & \\
\hline
2. Estrés Salino & Germinación ($\%$) & & & & & \\
(NaCl 150 mM) & Long. Radícula (cm) & & & & & \\
\hline
3. Oscuridad & Germinación ($\%$) & & & & & \\
(Agua Destilada) & Long. Radícula (cm) & & & & & \\
\hline
\end{tabular}%
}
\end{center}

\subsubsection{Tabla 2. Observaciones Morfológicas del Plántula y
Embrión}\label{tabla-2.-observaciones-morfoluxf3gicas-del-pluxe1ntula-y-embriuxf3n}

\begin{center}
\resizebox{\linewidth}{!}{%
\begin{tabular}{l|p{14em}}
\hline
\textbf{Tratamiento} & \textbf{Aspecto de la Radícula, Hipocótilo y Pigmentación} \\
\hline
1. Control (Luz) & \\
\hline
2. Estrés Salino & \\
\hline
3. Oscuridad & \\
\hline
\end{tabular}%
}
\end{center}

\subsection{\texorpdfstring{5. Análisis Estadístico y Gráfico (Estilo
\texttt{ggplot2})}{5. Análisis Estadístico y Gráfico (Estilo ggplot2)}}\label{anuxe1lisis-estaduxedstico-y-gruxe1fico-estilo-ggplot2}

Calcule el promedio (\(\bar{x}\)) y la desviación estándar (\(s\)) para
cada variable de respuesta:

\[\bar{x} = \frac{\sum_{i=1}^{n} x_i}{n} \qquad s = \sqrt{\frac{\sum_{i=1}^{n} (x_i - \bar{x})^2}{n - 1}}\]

\subsubsection{Gráfica 1: Porcentaje Final de Germinación (\%G) con
Barras de
Error}\label{gruxe1fica-1-porcentaje-final-de-germinaciuxf3n-g-con-barras-de-error}

\emph{Morfología:
\texttt{ggplot(data,\ aes(x=Tratamiento,\ y=Germinacion))\ +\ geom\_col()\ +\ geom\_errorbar()}}

\begin{center}
\begin{tikzpicture}[scale=0.85]
  % Grid ggplot2
  \draw[gray!20, step=0.5cm] (0,0) grid (6.5,4.0);
  \draw[gray!50, step=1cm] (0,0) grid (6.5,4.0);
  \draw[thick] (0,0) rectangle (6.5,4.0);
  
  % Eje X
  \node[below, font=\tiny] at (1.2,0) {\textbf{Control}};
  \node[below, font=\tiny] at (3.25,0) {\textbf{NaCl 150 mM}};
  \node[below, font=\tiny] at (5.3,0) {\textbf{Oscuridad}};
  \node[below=10pt, font=\small] at (3.25,0) {\textbf{Tratamiento Experimental}};
  
  % Eje Y
  \node[left, font=\tiny] at (0,0) {0};
  \node[left, font=\tiny] at (0,1) {25};
  \node[left, font=\tiny] at (0,2) {50};
  \node[left, font=\tiny] at (0,3) {75};
  \node[left, font=\tiny] at (0,4) {100};
  \node[rotate=90, above=10pt, font=\small] at (0,2.0) {\textbf{Germinación (\%)}};
\end{tikzpicture}
\end{center}

\subsubsection{Gráfica 2: Longitud Promedio de Radícula (cm) con Barras
de
Error}\label{gruxe1fica-2-longitud-promedio-de-raduxedcula-cm-con-barras-de-error}

\emph{Morfología:
\texttt{ggplot(data,\ aes(x=Tratamiento,\ y=Longitud))\ +\ geom\_col()\ +\ geom\_errorbar()}}

\begin{center}
\begin{tikzpicture}[scale=0.85]
  \draw[gray!20, step=0.5cm] (0,0) grid (6.5,4.0);
  \draw[gray!50, step=1cm] (0,0) grid (6.5,4.0);
  \draw[thick] (0,0) rectangle (6.5,4.0);
  
  \node[below, font=\tiny] at (1.2,0) {\textbf{Control}};
  \node[below, font=\tiny] at (3.25,0) {\textbf{NaCl 150 mM}};
  \node[below, font=\tiny] at (5.3,0) {\textbf{Oscuridad}};
  \node[below=10pt, font=\small] at (3.25,0) {\textbf{Tratamiento Experimental}};
  
  \node[left, font=\tiny] at (0,0) {0.0};
  \node[left, font=\tiny] at (0,1) {1.0};
  \node[left, font=\tiny] at (0,2) {2.0};
  \node[left, font=\tiny] at (0,3) {3.0};
  \node[left, font=\tiny] at (0,4) {4.0};
  \node[rotate=90, above=10pt, font=\small] at (0,2.0) {\textbf{Longitud Radícula (cm)}};
\end{tikzpicture}
\end{center}

\subsection{6. Actividades
Complementarias}\label{actividades-complementarias}

\begin{enumerate}
\def\labelenumi{\arabic{enumi}.}
\tightlist
\item
  Explique el mecanismo biofísico y metabólico por el cual un potencial
  osmótico negativo (\(\Psi_s\)) producido por el NaCl inhibe o retrasa
  la protrusión de la radícula.
\end{enumerate}

\hrulefill

\hrulefill

\begin{enumerate}
\def\labelenumi{\arabic{enumi}.}
\setcounter{enumi}{1}
\tightlist
\item
  ¿Qué papel cumple el fitocromo (forma \(P_{fr}\)) en la inducción de
  la germinación de semillas sensibles a la luz y cómo interactúa con el
  balance \(GA_3 / ABA\)?
\end{enumerate}

\hrulefill

\hrulefill

\begin{enumerate}
\def\labelenumi{\arabic{enumi}.}
\setcounter{enumi}{2}
\tightlist
\item
  Distinga entre germinación \textbf{epígea} e \textbf{hipógea} e
  identifique a cuál de estos tipos pertenece la especie vegetal
  utilizada en la práctica.
\end{enumerate}

\hrulefill

\hrulefill

\subsection{7. Conclusión}\label{conclusiuxf3n}

Interprete los resultados cuantitativos (\(\bar{x} \pm s\)) relacionando
la capacidad germinativa con la disponibilidad hídrica y las condiciones
de iluminación:

\hrulefill

\hrulefill

\hrulefill

\subsection{8. Bibliografía}\label{bibliografuxeda}

\begin{itemize}
\tightlist
\item
  Azcón-Bieto, J., \& Talón, M. (2008). \emph{Fundamentos de Fisiología
  Vegetal} (2.ª ed.). McGraw-Hill.
\item
  Bewley, J. D., Bradford, K. J., Hilhorst, H. W., \& Nonogaki, H.
  (2013). \emph{Seeds: Physiology of Development, Germination and
  Dormancy} (3rd ed.). Springer.
\item
  Taiz, L., \& Zeiger, E. (2002). \emph{Plant Physiology} (3rd ed.).
  Sinauer Associates.
\end{itemize}




\end{document}
